
\uline{où} $c$ \uline{est la
classe de Chern totale} (V 6.7).

Ce résultat se montre en recopiant la démonstration de 4.6, et en utilisant 5.6 (comparer [2], Exp 4).

\textbf{Corollaire 5.8.} \uline{Avec les notations et les hypothèses de} 5.1, \uline{on a pour} \uline{tout} $k$ :
\[
  f_*(\Fil^kK^{\bullet}(X))\subset \Fil^{k-r}K^{\bullet}(S).
\]

Conséquence immédiate de 5.3 et de la formule de projection.

\textbf{Proposition 5.9.} \uline{Avec les notations et les hypothèses de} 5.1, \uline{on a pour} \uline{tout} $i\ge 1$ :
\[
  f_*(\lambda^i(F))=0\qquad\bigl(\text{d'où }f_*(\gamma^i(F))=0\bigr).
\]
\uline{En particulier}, \uline{si on pose} $\lambda_{-1}(F)=\sum_{i=0}^{r}(-1)^i\lambda^i(F)$, \uline{on obtient} :
\[
  f_*(\lambda_{-1}(F))=1.
\]

On définit une application
\[
  \bar f_*:1+K^{\bullet}(X)[[t]]^+\longrightarrow 1+K^{\bullet}(S)[[t]]^+
\]
en posant :
\[
  \bar f_*\left(\sum_{i\ge0}a_it^i\right)=\sum_{i\ge0}f_*(a_i)t^i .
\]
Ce n'est pas en général un homomorphisme de groupes; néanmoins, la relation suivante est vérifiée pour $\varphi\in 1+K^{\bullet}(X)[[t]]^+$, $\psi\in 1+K^{\bullet}(S)[[t]]^+$ :
\[
  \bar f_*\bigl(\varphi.\hat G(f^*)(\psi)\bigr)=\psi.\bar f_*(\varphi).
\]
En effet, elle se vérifie coefficient par coefficient, et résulte immédiatement de la formule de projection.

De la relation $E=F+\xi$, on déduit donc :
